\documentclass[a4paper,12pt]{article}
\usepackage[utf8]{inputenc}
\usepackage[ngerman]{babel}
\usepackage{amsmath,amssymb}
\usepackage{geometry}\geometry{margin=2.5cm}
\usepackage{enumitem}
\usepackage{booktabs}
\usepackage{tikz}
\usepackage{pgfplots}\pgfplotsset{compat=1.18}
\usepackage{tcolorbox}
\tcbuselibrary{breakable}
\newtcolorbox{begriffe}[1][]{colback=yellow!10, colframe=yellow!60!black, fonttitle=\bfseries, title={Was bedeuten die Begriffe?}, breakable, #1}
\usepackage{xcolor}
\newcommand{\piv}[1]{{\color{red}\boxed{#1}}}

\title{\"Ubungsblatt 4 -- L\"osungen \\ \large Linear Algebra}
\author{}
\date{}

\begin{document}
\maketitle

%% ============================================================
\subsection*{Aufgabe 14}
%% ============================================================

\textbf{Angabe.} $A, B, C, M \in \mathbb{R}^{2 \times 2}$ mit
\[
A = \begin{pmatrix} 0 & 1 \\ 1 & 2 \end{pmatrix},\quad
B = \begin{pmatrix} 3 & 2 \\ -1 & 0 \end{pmatrix},\quad
C = \begin{pmatrix} 3 & 6 \\ 3 & 8 \end{pmatrix},\quad
M = \begin{pmatrix} 6 & 7 \\ 1 & 6 \end{pmatrix}.
\]
Ist $M \in \text{span}(A,B,C)$? Falls ja: eindeutig oder auf unendlich viele Arten?

\begin{begriffe}
\begin{itemize}[leftmargin=*, nosep]
    \item $\mathbb{R}^{2 \times 2}$ -- Menge aller $2 \times 2$-Matrizen mit reellen Eintr\"agen; bildet einen Vektorraum.
    \item $\text{span}(A,B,C)$ -- Menge aller Matrizen der Form $aA + bB + cC$ mit $a,b,c \in \mathbb{R}$.
    \item $M \in \text{span}(A,B,C)$ -- es existieren $a,b,c \in \mathbb{R}$ mit $aA + bB + cC = M$.
    \item \textbf{Eindeutige} Darstellung $=$ genau ein Tripel $(a,b,c)$; \textbf{unendlich viele} $=$ parametrische L\"osungsmenge.
\end{itemize}
\end{begriffe}

\medskip

\textbf{Schritt 1: Gleichungssystem aus $aA + bB + cC = M$ aufstellen.}

Eintragsweiser Vergleich liefert:
\[
\begin{cases}
3b + 3c = 6 & \text{(oben-links)} \\
a + 2b + 6c = 7 & \text{(oben-rechts)} \\
a - b + 3c = 1 & \text{(unten-links)} \\
2a + 8c = 6 & \text{(unten-rechts)}
\end{cases}
\]

\textbf{Schritt 2: Erweiterte Koeffizientenmatrix und Gau\ss-Elimination.}

\[
\left(\begin{array}{ccc|c}
0 & 3 & 3 & 6 \\
1 & 2 & 6 & 7 \\
1 & -1 & 3 & 1 \\
2 & 0 & 8 & 6
\end{array}\right)
\xrightarrow{Z_1 \leftrightarrow Z_2}
\left(\begin{array}{ccc|c}
1 & 2 & 6 & 7 \\
0 & 3 & 3 & 6 \\
1 & -1 & 3 & 1 \\
2 & 0 & 8 & 6
\end{array}\right)
\xrightarrow[Z_4 \to Z_4 - 2Z_1]{Z_3 \to Z_3 - Z_1}
\left(\begin{array}{ccc|c}
1 & 2 & 6 & 7 \\
0 & 3 & 3 & 6 \\
0 & -3 & -3 & -6 \\
0 & -4 & -4 & -8
\end{array}\right)
\]

\[
\xrightarrow[Z_4 \to 3Z_4 + 4Z_2]{Z_3 \to Z_3 + Z_2}
\left(\begin{array}{ccc|c}
1 & 2 & 6 & 7 \\
0 & 3 & 3 & 6 \\
0 & 0 & 0 & 0 \\
0 & 0 & 0 & 0
\end{array}\right)
\]

\textbf{Schritt 3: Pivots identifizieren und auswerten.}

In jeder Zeile der Stufenform ist der erste Eintrag $\neq 0$ ein \textbf{Pivot}. Nullzeilen liefern kein Pivot.

\[
\left(\begin{array}{ccc|c}
\piv{1} & 2 & 6 & 7 \\
0 & \piv{3} & 3 & 6 \\
0 & 0 & 0 & 0 \\
0 & 0 & 0 & 0
\end{array}\right)
\qquad
\begin{array}{l}
\text{Pivot-Spalten:}\; a, b \\
\text{Spalte ohne Pivot:}\; c \\
\Rightarrow c \;\text{frei, insgesamt 1 freie Variable}
\end{array}
\]

Keine Widerspruchszeile der Form $(0,\dots,0 \mid \neq 0)$ $\Rightarrow$ das System ist konsistent.
Mit einer freien Variable gibt es \textbf{unendlich viele L\"osungen}.

\medskip

\textbf{Schritt 4: Allgemeine L\"osung.}

Die Zeilen der Stufenform \"ubersetzen sich in folgende Gleichungen:
\begin{align*}
\text{Zeile 1:} &\quad 1 \cdot a + 2 \cdot b + 6 \cdot c = 7 \\
\text{Zeile 2:} &\quad 0 \cdot a + 3 \cdot b + 3 \cdot c = 6 \\
\text{Zeile 3:} &\quad 0 = 0 \quad \text{(keine Information)} \\
\text{Zeile 4:} &\quad 0 = 0 \quad \text{(keine Information)}
\end{align*}

Freie Variable setzen: $c = t$ mit $t \in \mathbb{R}$.

R\"uckw\"artseinsetzen von unten nach oben:
\begin{align*}
\text{Aus Zeile 2:} \quad 3b + 3t = 6 &\;\Rightarrow\; b = 2 - t \\
\text{Aus Zeile 1:} \quad a + 2(2-t) + 6t = 7 &\;\Rightarrow\; a + 4 + 4t = 7 \;\Rightarrow\; a = 3 - 4t
\end{align*}

\[
\fbox{\parbox{0.93\textwidth}{\centering
$M \in \text{span}(A,B,C)$ auf \textbf{unendlich viele Arten}:\\[2pt]
$(a,b,c) = (3-4t,\; 2-t,\; t),\quad t \in \mathbb{R}.$
}}
\]

\textbf{Schritt 5: Drei konkrete Linearkombinationen.}

Die Angabe verlangt drei konkrete Wege, $M$ als Linearkombination zu schreiben.
Vorgehen: Ich w\"ahle drei beliebige Werte f\"ur $t$, setze diese in die Formel
$(a,b,c) = (3-4t,\; 2-t,\; t)$ ein und berechne $aA + bB + cC$.
Als Kontrolle muss dabei jedes Mal $M$ herauskommen.

\begin{itemize}[leftmargin=*]
\item $t = 0 \Rightarrow (a,b,c) = (3, 2, 0)$:
\[
3A + 2B = \begin{pmatrix} 0 & 3 \\ 3 & 6 \end{pmatrix} + \begin{pmatrix} 6 & 4 \\ -2 & 0 \end{pmatrix}
= \begin{pmatrix} 6 & 7 \\ 1 & 6 \end{pmatrix} = M \;\checkmark
\]

\item $t = 1 \Rightarrow (a,b,c) = (-1, 1, 1)$:
\[
-A + B + C = \begin{pmatrix} 0 & -1 \\ -1 & -2 \end{pmatrix} + \begin{pmatrix} 3 & 2 \\ -1 & 0 \end{pmatrix} + \begin{pmatrix} 3 & 6 \\ 3 & 8 \end{pmatrix}
= \begin{pmatrix} 6 & 7 \\ 1 & 6 \end{pmatrix} = M \;\checkmark
\]

\item $t = -1 \Rightarrow (a,b,c) = (7, 3, -1)$:
\[
7A + 3B - C = \begin{pmatrix} 0 & 7 \\ 7 & 14 \end{pmatrix} + \begin{pmatrix} 9 & 6 \\ -3 & 0 \end{pmatrix} - \begin{pmatrix} 3 & 6 \\ 3 & 8 \end{pmatrix}
= \begin{pmatrix} 6 & 7 \\ 1 & 6 \end{pmatrix} = M \;\checkmark
\]
\end{itemize}

\bigskip
\hrule
\bigskip

%% ============================================================
\subsection*{Aufgabe 15(a) -- Polynome in $\mathbb{R}_2[x]$}
%% ============================================================

\textbf{Angabe.} Ist die Menge
\[
\{(x+1)(x-2),\; (x+1)(x-3),\; (x+1)(x-4)\}
\]
in $\mathbb{R}_2[x]$ linear unabh\"angig?

\begin{begriffe}
\begin{itemize}[leftmargin=*, nosep]
    \item $\mathbb{R}_2[x]$ -- Menge aller Polynome vom Grad $\leq 2$ mit reellen Koeffizienten.
    \item \textbf{Linear unabh\"angig} $\iff$ die Gleichung $\alpha p_1 + \beta p_2 + \gamma p_3 = 0$ (Nullpolynom) hat nur die triviale L\"osung $\alpha = \beta = \gamma = 0$.
    \item \textbf{Linear abh\"angig} $\iff$ es existiert eine \textbf{nichttriviale} Kombination (mindestens ein Koeffizient $\neq 0$), die das Nullpolynom ergibt.
    \item Gleichheit von Polynomen: zwei Polynome sind gleich $\iff$ alle Koeffizienten in der Basis $\{1, x, x^2\}$ \"ubereinstimmen (Koeffizientenvergleich).
\end{itemize}
\end{begriffe}

\medskip

\textbf{Schritt 1: Polynome ausmultiplizieren.}

Damit ich Koeffizientenvergleich machen kann, brauche ich die Polynome in
Standardform (sortiert nach Potenzen von $x$):
\begin{align*}
p_1 &= (x+1)(x-2) = x^2 - x - 2 \\
p_2 &= (x+1)(x-3) = x^2 - 2x - 3 \\
p_3 &= (x+1)(x-4) = x^2 - 3x - 4
\end{align*}

\textbf{Schritt 2: Ansatz $\alpha p_1 + \beta p_2 + \gamma p_3 = 0$ und Koeffizientenvergleich.}

Zwei Polynome sind gleich $\iff$ alle Koeffizienten zur gleichen Potenz \"ubereinstimmen.
Da die rechte Seite das Nullpolynom ist, muss jeder Koeffizient einzeln $0$ ergeben:
\[
\begin{cases}
\alpha + \beta + \gamma = 0 & \text{(Koeffizient von $x^2$)} \\
-\alpha - 2\beta - 3\gamma = 0 & \text{(Koeffizient von $x$)} \\
-2\alpha - 3\beta - 4\gamma = 0 & \text{(konstanter Term)}
\end{cases}
\]

\textbf{Schritt 3: Erweiterte Koeffizientenmatrix und Gau\ss-Elimination.}

\[
\left(\begin{array}{ccc|c}
1 & 1 & 1 & 0 \\
-1 & -2 & -3 & 0 \\
-2 & -3 & -4 & 0
\end{array}\right)
\xrightarrow[Z_3 \to Z_3 + 2Z_1]{Z_2 \to Z_2 + Z_1}
\left(\begin{array}{ccc|c}
1 & 1 & 1 & 0 \\
0 & -1 & -2 & 0 \\
0 & -1 & -2 & 0
\end{array}\right)
\xrightarrow{Z_3 \to Z_3 - Z_2}
\left(\begin{array}{ccc|c}
1 & 1 & 1 & 0 \\
0 & -1 & -2 & 0 \\
0 & 0 & 0 & 0
\end{array}\right)
\]

\textbf{Schritt 4: Pivots identifizieren und auswerten.}

\[
\left(\begin{array}{ccc|c}
\piv{1} & 1 & 1 & 0 \\
0 & \piv{-1} & -2 & 0 \\
0 & 0 & 0 & 0
\end{array}\right)
\qquad
\begin{array}{l}
\text{Pivot-Spalten:}\; \alpha, \beta \\
\text{Spalte ohne Pivot:}\; \gamma \\
\Rightarrow \gamma \;\text{frei}
\end{array}
\]

2 Pivots bei 3 Unbekannten $\Rightarrow$ es gibt eine freie Variable $\Rightarrow$ nichttriviale L\"osung existiert $\Rightarrow$ Menge ist \textbf{linear abh\"angig}.

\medskip

\textbf{Schritt 5: Konkrete nichttriviale Linearkombination bestimmen.}

Die Zeilen der Stufenform \"ubersetzen sich in:
\begin{align*}
\text{Zeile 1:} &\quad 1 \cdot \alpha + 1 \cdot \beta + 1 \cdot \gamma = 0 \\
\text{Zeile 2:} &\quad 0 \cdot \alpha + (-1) \cdot \beta + (-2) \cdot \gamma = 0 \\
\text{Zeile 3:} &\quad 0 = 0 \quad \text{(keine Information)}
\end{align*}

Die freie Variable darf ich frei w\"ahlen. Ich nehme den einfachsten Wert $\gamma = 1$
(damit ich nichttrivial bleibe):

R\"uckw\"artseinsetzen:
\begin{align*}
\text{Aus Zeile 2:} \quad -\beta - 2 \cdot 1 = 0 &\;\Rightarrow\; \beta = -2 \\
\text{Aus Zeile 1:} \quad \alpha + (-2) + 1 = 0 &\;\Rightarrow\; \alpha = 1
\end{align*}

\[
\fbox{\parbox{0.93\textwidth}{\centering
Menge ist \textbf{linear abh\"angig}:\\[2pt]
$1 \cdot p_1 - 2 \cdot p_2 + 1 \cdot p_3 = 0.$
}}
\]

\textbf{Schritt 6: Probe durch direktes Nachrechnen.}
\[
(x^2 - x - 2) - 2(x^2 - 2x - 3) + (x^2 - 3x - 4)
= (1-2+1)x^2 + (-1+4-3)x + (-2+6-4) = 0 \;\checkmark
\]

\bigskip
\hrule
\bigskip

%% ============================================================
\subsection*{Aufgabe 15(b) -- Funktionen $\{1, 2^x, 3^x\}$}
%% ============================================================

\textbf{Angabe.} Ist die Menge $\{1,\; 2^x,\; 3^x\}$ in $\text{FUNC}(\mathbb{R})$ linear unabh\"angig?

\begin{begriffe}
\begin{itemize}[leftmargin=*, nosep]
    \item $\text{FUNC}(\mathbb{R})$ -- Vektorraum aller Funktionen $f : \mathbb{R} \to \mathbb{R}$.
    \item $1$ bezeichnet die konstante Funktion $f(x) = 1$; $2^x$ und $3^x$ die Exponentialfunktionen.
    \item \textbf{Nullvektor} in $\text{FUNC}(\mathbb{R})$ $=$ Nullfunktion ($f(x) = 0$ f\"ur alle $x$).
    \item Gleichung $\alpha \cdot 1 + \beta \cdot 2^x + \gamma \cdot 3^x = 0$ muss f\"ur \textbf{alle} $x \in \mathbb{R}$ gleichzeitig gelten.
\end{itemize}
\end{begriffe}

\medskip

\textbf{Schritt 1: Ansatz und Umwandlung in ein l\"osbares System.}

Die Bedingung lautet:
\[
\alpha \cdot 1 + \beta \cdot 2^x + \gamma \cdot 3^x = 0 \quad \text{f\"ur \textbf{alle} } x \in \mathbb{R}.
\]

Eine Gleichung, die f\"ur alle $x$ gelten muss, gilt insbesondere f\"ur jede einzelne Wahl von $x$.
Ich setze also drei einfache Werte ein ($x = 0, 1, 2$) und erhalte drei
``normale'' lineare Gleichungen:

\begin{align*}
x = 0: &\quad \alpha \cdot 1 + \beta \cdot 2^0 + \gamma \cdot 3^0 = 0 \;\Rightarrow\; \alpha + \beta + \gamma = 0 \\
x = 1: &\quad \alpha \cdot 1 + \beta \cdot 2^1 + \gamma \cdot 3^1 = 0 \;\Rightarrow\; \alpha + 2\beta + 3\gamma = 0 \\
x = 2: &\quad \alpha \cdot 1 + \beta \cdot 2^2 + \gamma \cdot 3^2 = 0 \;\Rightarrow\; \alpha + 4\beta + 9\gamma = 0
\end{align*}

\textbf{Schritt 2: Erweiterte Koeffizientenmatrix und Gau\ss-Elimination.}

\[
\left(\begin{array}{ccc|c}
1 & 1 & 1 & 0 \\
1 & 2 & 3 & 0 \\
1 & 4 & 9 & 0
\end{array}\right)
\xrightarrow[Z_3 \to Z_3 - Z_1]{Z_2 \to Z_2 - Z_1}
\left(\begin{array}{ccc|c}
1 & 1 & 1 & 0 \\
0 & 1 & 2 & 0 \\
0 & 3 & 8 & 0
\end{array}\right)
\xrightarrow{Z_3 \to Z_3 - 3Z_2}
\left(\begin{array}{ccc|c}
1 & 1 & 1 & 0 \\
0 & 1 & 2 & 0 \\
0 & 0 & 2 & 0
\end{array}\right)
\]

\textbf{Schritt 3: Pivots identifizieren und auswerten.}

\[
\left(\begin{array}{ccc|c}
\piv{1} & 1 & 1 & 0 \\
0 & \piv{1} & 2 & 0 \\
0 & 0 & \piv{2} & 0
\end{array}\right)
\qquad
\begin{array}{l}
\text{Pivot-Spalten:}\; \alpha, \beta, \gamma \\
\text{Keine Spalte ohne Pivot} \\
\Rightarrow \text{keine freie Variable}
\end{array}
\]

3 Pivots bei 3 Unbekannten $\Rightarrow$ das System hat nur die triviale L\"osung $\alpha = \beta = \gamma = 0$.

\medskip

\textbf{Schritt 4: Folgerung f\"ur die urspr\"ungliche Bedingung.}

Ich habe gezeigt: Wenn die Gleichung bei $x = 0, 1, 2$ gilt, dann bleibt nur $\alpha = \beta = \gamma = 0$ \"ubrig.
Da die urspr\"ungliche Bedingung ``f\"ur alle $x$'' st\"arker ist (die drei Punkte sind nur ein Spezialfall davon),
kann erst recht keine nichttriviale L\"osung existieren.

\[
\fbox{\parbox{0.93\textwidth}{\centering
Menge ist \textbf{linear unabh\"angig}.
}}
\]

\bigskip
\hrule
\bigskip

%% ============================================================
\subsection*{Aufgabe 16}
%% ============================================================

\textbf{Angabe.} F\"ur welche $k \in \mathbb{R}$ ist die Menge
\[
\{v_1, v_2, v_3\} = \{(1,4,3,2),\; (1,2,3,4),\; (2, k+2, k-2, 2)\}
\]
in $\mathbb{R}^4$ linear unabh\"angig? Beide Methoden aus der Vorlesung zeigen.

\begin{begriffe}
\begin{itemize}[leftmargin=*, nosep]
    \item $\mathbb{R}^4$ -- Vektorraum aller 4-Tupel reeller Zahlen.
    \item \textbf{Parameter $k \in \mathbb{R}$} -- unbestimmte reelle Zahl; die Untersuchung erfolgt in Abh\"angigkeit von $k$ (Fallunterscheidung).
    \item \textbf{Beide Methoden aus der Vorlesung}:
    \begin{itemize}[nosep]
        \item Methode 1 -- Ansatz $\alpha v_1 + \beta v_2 + \gamma v_3 = 0$; L\"osen des homogenen Systems in $\alpha, \beta, \gamma$.
        \item Methode 2 -- Vektoren als Zeilen in eine Matrix schreiben; per Gau\ss{} auf Stufenform bringen und Rang ablesen.
    \end{itemize}
    \item \textbf{Pivot} -- f\"uhrender $\neq 0$-Eintrag einer Zeile in Stufenform. Linear unabh\"angig $\iff$ Anzahl der Pivots $=$ Anzahl der Vektoren.
\end{itemize}
\end{begriffe}

\bigskip

\subsubsection*{Methode 1: Homogene Gleichung $\alpha v_1 + \beta v_2 + \gamma v_3 = 0$}

\textbf{Schritt 1: Eintragsweise Bedingungen aufstellen.}

Die Gleichung $\alpha v_1 + \beta v_2 + \gamma v_3 = \vec 0$ bedeutet, dass jede der 4
Komponenten einzeln null ergibt. Durch Einsetzen der Vektoren erhalte ich:
\[
\begin{cases}
1\alpha + 1\beta + 2\gamma = 0 & (\text{Komponente 1: } 1, 1, 2) \\
4\alpha + 2\beta + (k+2)\gamma = 0 & (\text{Komponente 2: } 4, 2, k+2) \\
3\alpha + 3\beta + (k-2)\gamma = 0 & (\text{Komponente 3: } 3, 3, k-2) \\
2\alpha + 4\beta + 2\gamma = 0 & (\text{Komponente 4: } 2, 4, 2)
\end{cases}
\]

\textbf{Schritt 2: Gau\ss-Elimination mit Parameter $k$.}

Der Parameter $k$ bleibt stehen und wird mitgeschleppt -- nichts anderes macht diese Rechnung
schwieriger als ``normal''.

\[
\left(\begin{array}{ccc|c}
1 & 1 & 2 & 0 \\
4 & 2 & k+2 & 0 \\
3 & 3 & k-2 & 0 \\
2 & 4 & 2 & 0
\end{array}\right)
\xrightarrow[\substack{Z_3 \to Z_3 - 3Z_1 \\ Z_4 \to Z_4 - 2Z_1}]{Z_2 \to Z_2 - 4Z_1}
\left(\begin{array}{ccc|c}
1 & 1 & 2 & 0 \\
0 & -2 & k-6 & 0 \\
0 & 0 & k-8 & 0 \\
0 & 2 & -2 & 0
\end{array}\right)
\]

\[
\xrightarrow{Z_4 \to Z_4 + Z_2}
\left(\begin{array}{ccc|c}
1 & 1 & 2 & 0 \\
0 & -2 & k-6 & 0 \\
0 & 0 & k-8 & 0 \\
0 & 0 & k-8 & 0
\end{array}\right)
\xrightarrow{Z_4 \to Z_4 - Z_3}
\left(\begin{array}{ccc|c}
\piv{1} & 1 & 2 & 0 \\
0 & \piv{-2} & k-6 & 0 \\
0 & 0 & \piv{k-8} & 0 \\
0 & 0 & 0 & 0
\end{array}\right)
\]

\textbf{Schritt 3: Pivots identifizieren und Fallunterscheidung.}

Ich schaue in jeder Zeile auf den ersten Eintrag $\neq 0$ -- das sind die Pivot-Kandidaten:

\begin{itemize}[leftmargin=*, nosep]
\item Zeile 1: Pivot $= 1$ (immer $\neq 0$). \checkmark
\item Zeile 2: Pivot $= -2$ (immer $\neq 0$). \checkmark
\item Zeile 3: Pivot $= k-8$ -- dieser h\"angt von $k$ ab! Nur $\neq 0$, wenn $k \neq 8$.
\item Zeile 4: Nullzeile -- kein Pivot.
\end{itemize}

Fallunterscheidung nach dem Wert von $k$:

\begin{itemize}[leftmargin=*, nosep]
\item $k \neq 8$: drei Pivots, keine Spalte ohne Pivot $\Rightarrow$ nur triviale L\"osung $\Rightarrow$ linear \textbf{unabh\"angig}.
\item $k = 8$: nur zwei Pivots, Spalte $\gamma$ ohne Pivot $\Rightarrow$ eine freie Variable $\Rightarrow$ nichttriviale L\"osung existiert $\Rightarrow$ linear \textbf{abh\"angig}.
\end{itemize}

\[
\text{Linear unabh\"angig} \iff k - 8 \neq 0 \iff k \neq 8.
\]

\textbf{Schritt 4: Nichttriviale Kombination bei $k = 8$.}

F\"ur $k = 8$ ist $k - 8 = 0$ und $k - 6 = 2$. Die Zeilen der Stufenform ergeben:
\begin{align*}
\text{Zeile 1:} &\quad 1 \cdot \alpha + 1 \cdot \beta + 2 \cdot \gamma = 0 \\
\text{Zeile 2:} &\quad 0 \cdot \alpha + (-2) \cdot \beta + 2 \cdot \gamma = 0 \\
\text{Zeile 3:} &\quad 0 = 0 \quad \text{(keine Information, da $k-8 = 0$)} \\
\text{Zeile 4:} &\quad 0 = 0
\end{align*}

Freie Variable setzen: $\gamma = 1$.

R\"uckw\"artseinsetzen:
\begin{align*}
\text{Aus Zeile 2:} \quad -2\beta + 2 = 0 &\;\Rightarrow\; \beta = 1 \\
\text{Aus Zeile 1:} \quad \alpha + 1 + 2 = 0 &\;\Rightarrow\; \alpha = -3
\end{align*}
Probe: $-3(1,4,3,2) + (1,2,3,4) + (2,10,6,2) = (0,0,0,0) \;\checkmark$

\bigskip

\subsubsection*{Methode 2: Vektoren als Zeilen, Rang per Gau\ss{}}

\textbf{Schritt 1: Vektoren als Zeilen anordnen.}

Ich schreibe die drei Vektoren einfach untereinander. Die Idee: Wenn beim Gau\ss-Verfahren
eine der Zeilen zur Nullzeile wird, war der entsprechende Vektor eine Linearkombination
der anderen (also \"uberfl\"ussig) $\Rightarrow$ linear abh\"angig.

\[
\left(\begin{array}{cccc}
1 & 4 & 3 & 2 \\
1 & 2 & 3 & 4 \\
2 & k+2 & k-2 & 2
\end{array}\right)
\]

\textbf{Schritt 2: Gau\ss-Elimination.}

\[
\xrightarrow[Z_3 \to Z_3 - 2Z_1]{Z_2 \to Z_2 - Z_1}
\left(\begin{array}{cccc}
1 & 4 & 3 & 2 \\
0 & -2 & 0 & 2 \\
0 & k-6 & k-8 & -2
\end{array}\right)
\xrightarrow{Z_3 \to 2Z_3 + (k-6)Z_2}
\left(\begin{array}{cccc}
\piv{1} & 4 & 3 & 2 \\
0 & \piv{-2} & 0 & 2 \\
0 & 0 & \piv{2(k-8)} & 2(k-8)
\end{array}\right)
\]

Rechnung zu Zeile 3 (Schritt 2):
\begin{itemize}[nosep, leftmargin=*]
\item Spalte 2: $2(k-6) + (k-6) \cdot (-2) = 0$
\item Spalte 3: $2(k-8) + (k-6) \cdot 0 = 2(k-8)$
\item Spalte 4: $2 \cdot (-2) + (k-6) \cdot 2 = -4 + 2k - 12 = 2(k-8)$
\end{itemize}

\textbf{Schritt 3: Pivots identifizieren und auswerten.}

In jeder Zeile den ersten Eintrag $\neq 0$ suchen:
\begin{itemize}[leftmargin=*, nosep]
\item Zeile 1: Pivot $= 1$ (immer $\neq 0$). \checkmark
\item Zeile 2: Pivot $= -2$ (immer $\neq 0$). \checkmark
\item Zeile 3: Pivot $= 2(k-8)$ -- h\"angt von $k$ ab, wird $0$ genau bei $k = 8$.
\end{itemize}

Fallunterscheidung:

\begin{itemize}[leftmargin=*, nosep]
\item $k \neq 8$: drei Pivots, keine Nullzeile $\Rightarrow$ Rang $= 3 \Rightarrow$ linear \textbf{unabh\"angig}.
\item $k = 8$: Zeile 3 wird komplett zur Nullzeile $\Rightarrow$ Rang $= 2 < 3 \Rightarrow$ linear \textbf{abh\"angig}.
\end{itemize}

\bigskip

\[
\fbox{\parbox{0.93\textwidth}{\centering
Beide Methoden: $\{v_1, v_2, v_3\}$ ist \textbf{linear unabh\"angig} genau dann, wenn $k \neq 8$.\\[2pt]
Bei $k = 8$: $-3 v_1 + v_2 + v_3 = 0$ zeigt die Abh\"angigkeit.
}}
\]

\bigskip
\hrule
\bigskip

%% ============================================================
\subsection*{Aufgabe 17(a)}
%% ============================================================

\textbf{Angabe.} Sei $\{v_1, v_2, v_3\}$ linear unabh\"angig. Ist dann auch $\{v_1+v_2,\; v_2+v_3,\; v_2+2v_3\}$ linear unabh\"angig?

\begin{begriffe}
\begin{itemize}[leftmargin=*, nosep]
    \item $V$ -- beliebiger Vektorraum; $v_1, v_2, v_3 \in V$ beliebige Vektoren mit der angegebenen Voraussetzung.
    \item \textbf{Voraussetzung} ``$\{v_1, v_2, v_3\}$ linear unabh\"angig'' $\Rightarrow$ aus $a_1 v_1 + a_2 v_2 + a_3 v_3 = 0$ folgt $a_1 = a_2 = a_3 = 0$.
    \item \textbf{Zu zeigen:} gleiche Eigenschaft f\"ur die neuen Vektoren $v_1+v_2,\; v_2+v_3,\; v_2+2v_3$.
    \item \textbf{Strategie:} Linearkombination der neuen Vektoren nach $v_1, v_2, v_3$ umsortieren, dann Voraussetzung anwenden.
\end{itemize}
\end{begriffe}

\medskip

\textbf{Schritt 1: Ansatz.}

Um lineare Unabh\"angigkeit der neuen Menge zu zeigen, setze ich
eine beliebige Linearkombination der drei neuen Vektoren gleich null:
\[
\alpha(v_1+v_2) + \beta(v_2+v_3) + \gamma(v_2+2v_3) = 0.
\]
Ziel: Zeigen, dass hieraus zwingend $\alpha = \beta = \gamma = 0$ folgt.

\textbf{Schritt 2: Nach $v_1, v_2, v_3$ umsortieren.}

Ich multipliziere aus und sammle die Terme nach $v_1, v_2, v_3$. So steht vor jedem
``Basisvektor'' genau ein Ausdruck in $\alpha, \beta, \gamma$:
\[
\alpha v_1 + (\alpha + \beta + \gamma) v_2 + (\beta + 2\gamma) v_3 = 0
\]

\textbf{Schritt 3: Voraussetzung ausnutzen.}

Die Voraussetzung sagt: $\{v_1, v_2, v_3\}$ ist linear unabh\"angig. Das bedeutet:
eine Linearkombination dieser drei Vektoren ist nur dann null, wenn \textbf{jeder} Koeffizient
null ist. Damit erhalte ich ein konkretes Gleichungssystem in den Unbekannten $\alpha, \beta, \gamma$:
\[
\begin{cases}
\alpha = 0 & \text{(Koeffizient vor $v_1$)} \\
\alpha + \beta + \gamma = 0 & \text{(Koeffizient vor $v_2$)} \\
\beta + 2\gamma = 0 & \text{(Koeffizient vor $v_3$)}
\end{cases}
\]

\textbf{Schritt 4: System l\"osen.}

Das System ist einfach genug f\"ur Einsetzen von oben nach unten:
\begin{align*}
\text{Gleichung 1:} \quad \alpha &= 0 \\
\text{Gleichung 2 mit $\alpha = 0$:} \quad \beta + \gamma &= 0 \;\Rightarrow\; \beta = -\gamma \\
\text{Gleichung 3 mit $\beta = -\gamma$:} \quad -\gamma + 2\gamma = \gamma &= 0 \;\Rightarrow\; \beta = 0
\end{align*}

$\Rightarrow \alpha = \beta = \gamma = 0$. Nur triviale L\"osung.

\[
\fbox{\parbox{0.93\textwidth}{\centering
Aussage (a) ist \textbf{wahr}: $\{v_1+v_2,\; v_2+v_3,\; v_2+2v_3\}$ ist linear unabh\"angig.
}}
\]

\bigskip
\hrule
\bigskip

%% ============================================================
\subsection*{Aufgabe 17(b)}
%% ============================================================

\textbf{Angabe.} Ist $\{7v_1 - 11v_2,\; 11v_2 - 13v_3,\; 7v_1 - 13v_3\}$ notwendig linear abh\"angig?

\begin{begriffe}
\begin{itemize}[leftmargin=*, nosep]
    \item $V$ -- beliebiger Vektorraum; $v_1, v_2, v_3 \in V$ beliebig (keine Voraussetzung \"uber Abh\"angigkeit).
    \item \textbf{Notwendigerweise} linear abh\"angig $\iff$ die Menge ist linear abh\"angig \textbf{f\"ur jede} Wahl von $v_1, v_2, v_3$ (nicht nur f\"ur ein Beispiel).
    \item \textbf{Nachweis} durch Angabe einer konkreten, nichttrivialen Linearkombination der Null, deren Koeffizienten nicht von den $v_i$ abh\"angen.
\end{itemize}
\end{begriffe}

\medskip

\textbf{Schritt 1: Nach einem Zusammenhang zwischen den drei Vektoren suchen.}

Um ``notwendigerweise linear abh\"angig'' zu zeigen, reicht es, eine feste
Linearkombination der Null anzugeben, die f\"ur \textbf{jede} Wahl von $v_1, v_2, v_3$ gilt.
Ich probiere die einfachste: die Summe der ersten beiden Vektoren:
\[
(7v_1 - 11v_2) + (11v_2 - 13v_3) = 7v_1 \underbrace{-11v_2 + 11v_2}_{=0} - 13v_3 = 7v_1 - 13v_3.
\]

Die $v_2$-Terme haben sich aufgehoben. Das Ergebnis $7v_1 - 13v_3$ ist genau der dritte Vektor.

\medskip

\textbf{Schritt 2: Diese Beobachtung als Linearkombination der Null hinschreiben.}

``Vektor 1 $+$ Vektor 2 $=$ Vektor 3'' l\"asst sich umstellen zu
``Vektor 1 $+$ Vektor 2 $-$ Vektor 3 $= 0$'':
\[
1 \cdot (7v_1 - 11v_2) + 1 \cdot (11v_2 - 13v_3) + (-1) \cdot (7v_1 - 13v_3) = 0
\]

Koeffizienten $(1, 1, -1)$ -- alle ungleich null $\Rightarrow$ nichttriviale Kombination.

\medskip

\textbf{Schritt 3: Folgerung.}

Die Koeffizienten $(1, 1, -1)$ h\"angen \textbf{nicht} von den konkreten Werten von $v_1, v_2, v_3$ ab.
Die Gleichung ist f\"ur jede Wahl der drei Vektoren im gleichen Vektorraum erf\"ullt.

\[
\fbox{\parbox{0.93\textwidth}{\centering
Aussage (b) ist \textbf{wahr}: Die Menge ist f\"ur jede Wahl von $v_1, v_2, v_3$ linear abh\"angig.
}}
\]

\end{document}
